\chapter{Hướng dẫn tái lập}
\label{app:tai-lap}

\section{Cài đặt}
Từ gốc repo, chạy \repoFile{uv sync --group dev --extra analysis} để chuẩn bị môi trường.
Chạy \repoFile{uv run pytest} để kiểm tra. Nếu cần backend native, cài nhóm phụ thuộc
\repoFile{native} và build theo \repoFile{setup_native.py}.

\section{Chạy giải đấu và thí nghiệm}
Notebook \repoFile{notebooks/arena_analysis.ipynb} đọc giải cũ.
\par
Notebook \repoFile{notebooks/l7_l8_duel.ipynb} đọc log L7--L8 và có tùy chọn chạy arena đủ tám cấp.
Với từng lượt đấu mới, lưu seed, cấu hình, khai cuộc, backend và checksum checkpoint.

\section{Huấn luyện}
Bộ sinh chuyên gia, học có giám sát, REINFORCE và AlphaZero lần lượt nằm ở
\repoFile{caro/training/gen_expert_data.py}, \repoFile{train_sl.py},
\repoFile{train_pg.py} và \repoFile{alphazero_loop.py}. Lệnh chạy đầy đủ L8:
\begin{lstlisting}[language=bash]
uv run python -m caro.training.alphazero_loop training/alphazero
\end{lstlisting}
Dùng config \repoFile{training/alphazero_smoke} để kiểm tra đường chạy nhanh.

\section{Sinh bảng và hình}
Từ gốc repo, chạy \repoFile{python report/scripts/sinh-ho-so-caro.py} với môi trường của dự án
để sinh lại bảng và hình từ log gốc. Từ thư mục \repoFile{report/}, chạy
\repoFile{latexmk main.tex}. Mục lục, danh mục hình, danh mục bảng và trích dẫn vẫn theo template.

\section{Demo}
Từ gốc repo, chạy \repoFile{uv run uvicorn caro.ui.server:app}. Giao diện có hai chế độ
đối đầu và phát lại. Tác tử neural chỉ sẵn sàng khi nạp đúng checkpoint.

\section{Chỉ mục nguồn thực nghiệm}
\begin{itemize}
\item A1, A2, A7: các \repoFile{summary.json} và \repoFile{metrics.jsonl} dưới \repoFile{runs/bench/}.
\item A3, A4: log ván dưới \repoFile{runs/tournament/a3_mcts_sims/} và \repoFile{runs/tournament/a4_rollout/}.
\item Arena cũ: \repoFile{runs/arena_matrix/games.csv}; tổng hợp theo tác tử trong \repoFile{runs/arena_matrix/ranking.csv}.
\item L5: \repoFile{runs/train/sl_policy_v3/metrics.json}.
Đối chiếu checkpoint: \repoFile{runs/train_eval/results.csv}.
\item L6: \repoFile{runs/train/pg_policy/20261118-125103_v3_init/}.
\item A5: \repoFile{experiments/a5_puct_eval.py} và \repoFile{runs/eval/a5_puct/results.csv}.
\item L8: \repoFile{runs/train/alphazero/20261124-160615_gpu-native/}, gồm cấu hình, chỉ số qua vòng, tổng kết và ván tự đấu/kiểm tra.
\item L8--L3: \repoFile{runs/eval/tournament/l8_vs_l3/20261125-134954_round37/}.
\item Replay chiến thuật: \repoFile{notebooks/l7_l8_duel.ipynb}; nguồn luật nằm dưới \repoFile{caro/engine/}.
\end{itemize}

\section{Đối chiếu số liệu với log}
Khi kiểm chứng một con số, trước hết xác định mã lượt chạy và tham số trong cấu hình đã lưu.
Với log JSONL, lọc đúng cặp tác tử hoặc vòng, chuyển kết quả tuyệt đối của X về góc nhìn
tác tử được báo cáo, rồi cộng thắng, hòa và thua. Kiểm tra có đủ hai màu cho từng khai cuộc
trước khi tính bootstrap theo cặp. Đối với thời gian mỗi nước, lấy bên thực sự đi sau chuỗi
khai cuộc để chia các giá trị xen kẽ cho đúng tác tử.

Các ván minh họa trong Chương~\ref{ch:thuc-nghiem} có mã ván và mã khai cuộc.
Có thể dựng lại bằng \repoFile{CaroGame.from_moves} hoặc chọn đúng lượt trong
\repoFile{notebooks/l7_l8_duel.ipynb}. Với log tự đấu, mã ván có thể lặp giữa các nhóm,
nên đối chiếu đồng thời trường \repoFile{tournament_id} và \repoFile{game_id}.
